\section{The quadratic coefficient-prime character in the cyclotomic variable}
\label{sec:cubic-ramified-cyclotomic}

The character in this section is the $\ell$-primary
\emph{conductor component} of a Dirichlet character, not its
$\ell$-power-order component.  The actual source defines it
as \texttt{comp l psi} in
\texttt{DirichletPrimeToLPartMulLPowerPart.lean}; the condition
\texttt{GD.ine} in \texttt{Counts.lean} is its square being one.
These definitions are used below without a Lean execution.

\begin{lemma}[AC-K29-good-twist: the good ordinary quadratic twist]
\label{AC-K29-good-twist}
Let $\ell\ge5$ split in $K=\mathbb Q(\omega)$, let $e\ge1$,
and let $f$ be a cuspidal newform of weight $2e+1$ and
nebentypus $\eta_K$.  Put $V=M_f(e)$ in the cohomological
convention of (CER.2).  Suppose that $V|_{G_{\mathbb Q_\ell}}$
is potentially crystalline and has a stable line with
characters $\delta,\delta^{-1}$, in this order, where
\[
       \delta|_{I_\ell}=\tau|_{I_\ell}\kappa^e|_{I_\ell}.
                                                        \tag{CRC.1}
\]
Here $\kappa$ is the cyclotomic character and $\tau$ is a
Dirichlet character of conductor a power of $\ell$ with
$\tau^2=1$.  Let $f_0$ be the primitive newform associated
to the automorphic twist $f\otimes\tau$.  Then
\[
 M_{f_0}=M_f\tau,\qquad V_0:=M_{f_0}(e)=V\tau,
 \qquad \operatorname{neb}(f_0)=\eta_K.                    \tag{CRC.2}
\]
The level of $f_0$ is prime to $\ell$, and $f_0$ is
ordinary at the chosen coefficient prime.  Its two Hecke
roots have valuations $0$ and $2e$.
If $f$ is non-CM, so is $f_0$; their restrictions to every
derived Galois subgroup have the same image.
The character $\tau$ is either trivial or the quadratic
character of conductor exactly $\ell$.

In particular the lemma applies to an actual ordinary
Type-II point with the flag of (G48.8), potential
crystallinity as in that local condition, and
$\tau=\operatorname{comp}_\ell\psi$ satisfying
\texttt{GD.ine}.  No assertion that all points of an
unidentified component are ordinary is needed here.
\end{lemma}
\begin{proof}
The CRT idempotent in the displayed source definition
restricts a Dirichlet character to the $\ell$-power factor
of its modulus.  Thus its associated global character is
unramified away from $\ell$, and its restriction to
$I_\ell$ is that of $\psi$.  A character of order dividing
two on $\mathbb Z_\ell^\times$ is trivial on the pro-$\ell$
principal units.  If nontrivial, it is the unique quadratic
character of $\mathbb F_\ell^\times$, proving the conductor
assertion.  This also verifies directly the interpretation
of \texttt{GD.ine} used here.

The characters on the two lines of $V_0=V\tau$ are
$\delta\tau$ and $\delta^{-1}\tau$.  By (CRC.1) they
are respectively an unramified character times
$\kappa^e$ and its inverse.  Both are crystalline.
We verify that their possibly nonsplit extension is
crystalline as well.  Its class belongs to
$H^1(\mathbb Q_\ell,B)$, where $B$ is their crystalline
Hom character.  Choose a finite Galois extension $F$ over
which the extension is crystalline.  The image of this
class in $H^1(\mathbb Q_\ell,B_{\rm cris}\otimes B)$
restricts to zero over $F$.
Restriction on this latter group is injective: the
restriction--corestriction composite is multiplication
by $[F:\mathbb Q_\ell]$, which is invertible in the
coefficient field.  Equivalently use finite-group averaging
in the inflation--restriction sequence.
The extension class therefore lies in $H^1_f$ already
over $\mathbb Q_\ell$.  This is the defining crystalline
extension criterion, and proves that $V_0$ is crystalline.

Automorphic twisting preserves cuspidality and the
archimedean weight, changes the nebentypus by $\tau^2$,
and gives the Galois representation in (CRC.2), as follows
also from its eigenvalues away from the finite bad set
and Chebotarev.  Local--global compatibility at $\ell$
identifies the Frobenius-semisimplified Weil--Deligne
representation of $M_{f_0}$ with that of its local
newform representation.  The precise classical statement is Saito's
\emph{Modular forms and $p$-adic Hodge theory},
\emph{Inventiones}~129 (1997), 607--620, as explicitly
restated by its author in
\emph{Galois representations in arithmetic geometry},
\S2c, Theorem~3, printed p.~9 of
\url{https://www.ms.u-tokyo.ac.jp/~t-saito/pp/suams.pdf}.
That statement and its preceding paragraph concern
$V_{\ell,f}$ of $G_{\mathbb Q}$ and every prime $p$;
they identify $D_{\rm pst}(V_{f,p})^{F\text{-}\rm ss}$
with the local Langlands parameter of $\pi_{f,p}$.
Since a crystalline
representation has unramified Weil--Deligne parameter
with zero monodromy, the local representation is
spherical.  Thus the primitive level is prime to $\ell$.
This uses the stated compatibility theorem, not a new
proof of that theorem.

The first line of $M_{f_0}=V_0(-e)$ is unramified.
Its Frobenius eigenvalue is a unit, since a continuous
unramified character of the compact local Galois group
has compact image.  Its crystalline Frobenius slope is
therefore zero.  The determinant has the other slope
$2e$, giving the stated Hecke-root valuations and
ordinarity.  Quadratic scalar twisting preserves the
property of being induced from a quadratic field;
hence it preserves non-CM.  Finally every scalar
character is trivial on a derived subgroup, which
proves the last image assertion.
\end{proof}

\begin{theorem}[AC-K29-ramified: the original coefficients after cyclotomic specialization]
\label{AC-K29-ramified}
Retain AC-K29-good-twist.  Suppose $f$ is non-CM and has
the characteristic-zero derived-image conclusion of
AC-K29-split.  Suppose further that $\ell\ne269$ and
\[
 \chi(\operatorname{Fr}_{\mathfrak p})\ne1,
 \qquad \mathfrak p\mid\ell,                              \tag{CRC.3}
\]
so that the actual two CM ordinary quotients and their
integral ray norm maps are those of AC-K29-ray-class
and AC-K29-ray.  Then
\[
                  H^1_f(K,V\chi^a)=0\quad(a=1,2).          \tag{CRC.4}
\]
In particular the requirement that the original $f$
have level prime to $\ell$ can be replaced by the
actual potential-crystalline ordinary condition
(CRC.1).  The conclusion is over the coefficient field;
it supplies no uniform integral Euler-index bound.

More precisely there are two ray Euler systems for
\[
 U_a=(V|_{G_K}\chi^a)^*(1)
       =M_f^*|_{G_K}\chi^{-a}(1-e),                       \tag{CRC.5}
\]
with one fixed denominator independent of the auxiliary
ray modulus, Bloch--Kato local conditions away from
$\ell$, and exactly the norm polynomial
\[
 P_{\mathfrak q}(X)=
       \det(1-X\operatorname{Fr}_{\mathfrak q}^{-1}
                    \mid V|_{G_K}\chi^a).                \tag{CRC.6}
\]
Their two rational base classes span the singular local
space at $\ell$ after Shapiro.  The proof keeps the
ordinary CM form $g=\theta_\chi$ unchanged.
\end{theorem}
\begin{proof}
For $\tau=1$, AC-K29-good-twist gives a good ordinary
$f$, and AC-K29-primitive proves every assertion.  We
therefore treat the nontrivial character $\tau$, of
conductor $\ell$.

\proofstep{1. Specialize the actual ray and cyclotomic system.}
Construct the fixed-form system of AC-K29-ray for
$(f_0,g)$, retaining its cyclotomic Iwasawa variable
until the last step.  All of the integral CM ray maps
in that proof are compatible with this variable.
Let $j=e-1$ and specialize at the locally algebraic
character $\kappa^j\tau$.  The coefficient map is
\[
 M_{f_0}^*\otimes M_g^*(-j)\otimes\tau^{-1}
       =M_f^*\otimes M_g^*(1-e),                         \tag{CRC.7}
\]
using $\tau^{-1}=\tau$ and (CRC.2).
Under Shapiro this is exactly (CRC.5).
One can perform the finite part explicitly by taking
the $\tau$ idempotent at $\mu_\ell$ after the
algebraic cyclotomic twist; equivalently, before that
twist use the idempotent for
$\omega_{\rm Teich}^{\,j}\tau$.  Its denominator divides $\ell-1$
and is an $\ell$-adic unit.  The $\ell$-primary ray
fields of conductor prime to $\ell$ are disjoint from
this tame cyclotomic extension.  Thus this operation
adds no denominator depending on the ray modulus.

The norm identities hold before specialization.
On their group-ring Frobenius variable the finite
specialization inserts $\tau(\operatorname{Fr}_{\mathfrak q})$.
Consequently the polynomial for $V_0\chi^a$ becomes
that for $V_0\chi^a\tau=V\chi^a$, which is (CRC.6).
Corestriction and the fixed finite idempotent commute.
The fixed tame operators of AC-K29-primitive are
applied before the actual compatible CM quotients,
again as $\pi_{\mathfrak n}\circ T$; no commutation
with an unspecified quotient is asserted.
They commute with the finite cyclotomic projection
and with every new auxiliary-prime norm map because
their support is fixed and disjoint from those primes.
All denominator bounds from AC-K29-ray therefore
remain uniform in $\mathfrak n$.
The classes are still specializations of cyclotomic
universal norms, which gives the same away-$\ell$
Bloch--Kato local condition as in that proposition.

\proofstep{2. Primitive scalar reciprocity at the finite character.}
The finite tame-data proof of AC-K29-primitive-reciprocity
works at this character with the following precise
changes.  Its scalar identity is the ordinary KLZ
identity before cyclotomic specialization, so the
coefficient map in (CRC.7) is used on both sides.
In its second-slice argument (CEP.7)--(CEP.8), take
\[
 (\kappa_1,\kappa_2,\sigma)=(2e+1,1,e+\tau),
 \qquad \tau'=e-1+\tau .                                 \tag{CRC.8}
\]
The symbols in (CRC.8) denote locally algebraic
characters in additive weight notation.  The
near-overconvergence degree is still $e-1$;
$2e+1>2(e-1)$ excludes every projector pole.
Loeffler, arXiv:1704.04049v2, Lemma~5.3 and
Corollary~6.2 use locally algebraic slice characters,
including their finite components.  Thus exactly
the scalar two-slice identity used in the primitive
proof applies.  This is not a comparison between
ordinary and Coleman cohomology classes.

For the standard local data at $\ell$, the explicit
nontrivial-character factor is given by the same
paper's Proposition~2.10 and Appendix~A, Steps~1--4:
\[
 \mathcal E_\ell(f_0,g,e+\tau)
    =G(\tau)^2\frac{\ell^{2e-2}}
                    {\alpha_{f_0}^{2}\alpha_g\beta_g}
       \ne0.                                             \tag{CRC.9}
\]
The two denominator factors are
$1-\beta_{f_0}/(\ell\alpha_{f_0})$ and
$1-\beta_{f_0}/\alpha_{f_0}$; their subtracted terms
have valuations $2e-1$ and $2e$, so both are nonzero.
No local primitive Euler factor at $\ell$ is missing:
all four unramified tensor characters have acquired
the nontrivial scalar inertia $\tau$, and hence the
local invariant space is zero.

At the finitely many tame bad primes choose the
normalized nonzero Rankin data for
$\pi_{f_0}\otimes\pi_g\otimes\tau^{-1}$ as in
AC-K29-primitive-data.  Its central character is still
one since $\tau^2=1$.  At an auxiliary good prime
the same $\mathcal S_{00}$ quotient argument applies;
non-CM and the different Hodge--Tate weights give
an auxiliary prime with the required inequivalent
local factors.  All the normalized local
functional-equation factors remain nonzero.
These are finitely many fixed tame operators, so the
actual class construction of Step~1 supplies them.
The primitive global value is exactly
\[
 L^{\rm prim}(f_0,g,\tau^{-1},e)
           =L^{\rm prim}(f,g,e)\ne0.                     \tag{CRC.10}
\]
The nonvanishing is the primitive boundary argument
of AC-K29-bad; it has no hypothesis on the
coefficient-prime level of $f$.

Choose one permitted smoothing integer $c$ equal to
one modulo all fixed tame levels and with
$c^2\not\equiv1\pmod\ell$.
The Chinese remainder theorem permits this for
$\ell\ge5$.  The cyclotomic finite factor in the
smoothing expression is $\tau(c)^2=1$; its algebraic
exponent cancels at $j=e-1$.  The exact factor is
therefore $c^2-1\ne0$ on both rational base classes.
At higher ray modulus retain the original group-ring
smoothing operator; the same fixed $c$ is used and
its universal norm compatibility is unchanged.
Equations (CRC.9)--(CRC.10), the chosen tame periods,
and the ordinary scalar reciprocity identity prove
that both actual quotient regulators are nonzero.

\proofstep{3. The ramified quotient regulator and singular space.}
This specialization is expressly covered by
\emph{ordinary} KLZ, arXiv:1503.02888v6,
Theorem~8.2.8, printed pp.~43--44, whose proof
applies the functorial unramified construction of
Theorem~8.2.3 and then twists the cyclotomic variable.
The finite character there is
$\varepsilon_{g,\ell}\tau^{-1}=\tau$ with conductor
exponent one.  At $k'=-1$, $j=e-1$, its map is
\[
   (e-1)!\frac{\ell^e}{\alpha_{f_0}\beta_g}
             G(\tau)^{-1}\exp^*.                         \tag{CRC.11}
\]
The Gauss-period identification of the de Rham
ramified character with the unramified rank-one
module is part of that formula.  Every scalar in
(CRC.11) is nonzero.  In particular neither of the
unramified exceptional Euler factors is present.
The CM weight-one quotient and its functional are
the actual ones of AC-K29-ray and
AC-K29-primitive; no negative algebraic moment or
vanishing coefficient-module comparison is used.

Here are the local dimension and projection checks.
Write $\rho=\operatorname{Ind}_K^{\mathbb Q}\chi$,
$W=V\otimes\rho$, and $U=W^*(1)$.  Both $U$ and
$U^*(1)$ are crystalline representations twisted
by the nontrivial tame quadratic $\tau$.
Their potentially crystalline modules have scalar
inertia $\tau$, so
\[
 D_{\rm cris}(U)=D_{\rm cris}(U^*(1))=0,
 \quad H^0(\mathbb Q_\ell,U)=H^2(\mathbb Q_\ell,U)=0.
                                                        \tag{CRC.12}
\]
For example, pass to the totally ramified quadratic
extension killing $\tau$: inertia acts trivially on
its maximal unramified coefficient field and by
$\tau$ on the crystalline module, so its invariants
vanish.  Local duality gives the $H^2$ assertion.
The local Euler characteristic and Bloch--Kato
exact sequences now give
\[
 \begin{gathered}
 \dim H^1(\mathbb Q_\ell,U)=4,\qquad
 H^1_e=H^1_f=H^1_g,\qquad \dim H^1_f=2,
 \\
 \exp^*:H^1_s(\mathbb Q_\ell,U)
          \xrightarrow{\ \sim\ }\operatorname{Fil}^0D_{\rm dR}(U).
                                                        \tag{CRC.13}
 \end{gathered}
\]
These dimensions can also be read from the actual
ordinary exact sequence.  Its kernel in $U$ has
cyclotomic exponent $e+1$, and its rank-two quotient
has exponent $1-e$.  Finite characters have Hodge--Tate
weight zero.  Thus the kernel has zero
$\operatorname{Fil}^0$, the quotient has full
$\operatorname{Fil}^0$, and strictness of the de Rham
filtration makes projection an isomorphism on the
two-dimensional space in (CRC.13).

KLZ Proposition~8.1.7 annihilates the $F^{--}$
projection before any cyclotomic specialization.
The two CM refinements therefore put the two
singular images in the two different CM lines of
this rank-two quotient.  Formula (CRC.11) and
Step~2 make each image nonzero.  They span the
whole singular local space, with no unramified
$D_{\rm cris}$ substitution for the ramified
space in (CRC.13).

\proofstep{4. The strict Selmer conclusion.}
Apply the local Poitou--Tate argument (CER.7) to
these actual two classes.  It forces the localization
of $H^1_f(\mathbb Q,W)$ at $\ell$ to be zero; Shapiro
gives the strict-at-both-primes identity (CER.8) for
the original $R_a=V\chi^a$.
At least one base Euler class is nonzero by its
nonzero singular localization.  The rank-two
$K$-Euler-system bound used in AC-K29-lower then
annihilates this strict group: (CRC.5)--(CRC.6)
supply its actual coefficients and norm polynomials,
Step~1 supplies its fixed lattice denominator and
away-$\ell$ local conditions, and scalar twisting
leaves the derived-image hypotheses unchanged.
All ray and cyclotomic fields are abelian, so the
required derived-image elements act trivially on
them.  This proves (CRC.4).
\end{proof}

\begin{remark}[Scope and geometric dictionary]
The quadratic twist is tensoring an $\ell$-adic local
system by a rank-one finite local system; the
cyclotomic character projection cancels that tensor
factor on the actual cohomology classes.  Shapiro
uses the corresponding finite etale pushforward.
This is the usual representation--local-system
category dictionary, not a new attribution of the
Euler or regulator argument to Grothendieck.
The theorem retains (CRC.3), the actual ordinary
flag and the derived-image condition.  It does not
settle the split repeated-root CM case, the primes
$2,3,269$, integral index bounds, or any unspecified
Type-I automorphy input.
\end{remark}
